\documentclass[conference]{IEEEtran}
\IEEEoverridecommandlockouts

\usepackage[T1]{fontenc}
\usepackage[utf8]{inputenc}
\usepackage{microtype}
\usepackage{amsmath,amssymb}
\usepackage{booktabs}
\usepackage{array}
\usepackage{xurl}
\usepackage{url}
\usepackage{hyperref}
\hypersetup{hidelinks}
\urlstyle{same}

\begin{document}

\title{Improving SGFuzz State Instrumentation Using Clang AST Analysis}

\author{
\IEEEauthorblockN{Will D'Andrade}
\IEEEauthorblockA{Department of Cybersecurity\\
Rochester Institute of Technology\\
Rochester, New York, USA}
\and
\IEEEauthorblockN{Jose Luis Gonzalez Paniagua}
\IEEEauthorblockA{Department of Cybersecurity\\
Rochester Institute of Technology\\
Rochester, New York, USA}
}
\maketitle

\begin{abstract}
SGFuzz obtains stateful feedback by instrumenting assignments to state-related
enum variables. Its source pass uses regular expressions, which can miss modern
C++ representations and select unrelated assignments. We replace only this
selection step with a Clang abstract syntax tree (AST) analysis while retaining
the original SGFuzz callback and State Transition Tree (STT) runtime. The
prototype identifies project enum storage, resolves assignment targets to their
declarations, and rewrites accepted standalone assignments in a copied source
tree. In IXWebSocket, the regex selector placed zero callbacks at the target
atomic enum field; the AST selector placed five. With the same 36-unit build,
fixed corpus, and deterministic replay, the AST condition produced seven STT
nodes and two leaves, versus the control's root-only one node and one leaf.
Five pre-registered C and C++ regression controls passed, and static audits
parsed all 36 selected IXWebSocket C++ units and all 124 selected H2O C units
without failures. A completed 60-second sanitizer preflight sustained the
instrumented IXWebSocket harness at about 96k executions per second with no
sanitizer finding. These results support improved source-selection fidelity and
end-to-end feasibility; they do not establish superior fuzzing throughput,
coverage, or bug discovery.
\end{abstract}

\begin{IEEEkeywords}
Stateful Greybox Fuzzing, SGFuzz, Clang, Abstract Syntax Tree, Source
Instrumentation, Software Testing
\end{IEEEkeywords}

\section{Introduction}

Coverage-guided fuzzers use observations from prior executions to retain and
mutate inputs. For stateful programs, code coverage alone may not distinguish
executions that reach different protocol states through similar control flow.
SGFuzz addresses this problem by observing assignments to variables that
represent states and by storing the observed transition sequences in a State
Transition Tree (STT)~\cite{ba2022}.

The usefulness of this feedback depends on selecting the right source
statements. The original SGFuzz source pass searches text with regular
expressions. Text alone does not reliably identify the declared type or target
of an assignment. In particular, a scoped enum or \texttt{std::atomic<Enum>}
field can evade a text pattern, while an unrelated assignment with a similar
shape can be selected incorrectly. A missed write prevents the runtime from
observing a state transition; an incorrect rewrite can prevent a target from
building.

This paper evaluates a narrow replacement for the selection step. The prototype
uses Clang's AST to associate an assignment left-hand side with a declaration
and type, then emits the existing SGFuzz callback only for supported writes to
project-owned enum storage. It neither replaces SGFuzz's runtime nor claims to
infer the semantic importance of every enum. Our question is whether typed AST
analysis improves instrumentation fidelity without regressing cases that the
regex selector already handles.

\section{Background and Design}

\subsection{SGFuzz feedback}

SGFuzz inserts \texttt{\_\_sfuzzer\_instrument(location, value)} at selected
writes. The unchanged runtime receives those events and builds an STT. An STT
node count includes the root. A leaf is a distinct root-to-leaf transition
sequence, not an individual enum value, code edge, state declaration, or bug.
Thus, a root-only result of one node and one leaf means that no instrumented
state transition was observed by that execution.

The baseline paper reports stateful fuzzing results across protocol
implementations~\cite{ba2022}. This work does not repeat those performance or
bug-discovery claims. It isolates the earlier source-selection stage that
determines which state events reach the same callback ABI and runtime.

\subsection{AST-aware source selection}

The implementation has a detector and a copy-only instrumenter. The detector
invokes Clang with a textual AST dump for selected C or C++ translation units.
It records project enum declarations, then records \texttt{VarDecl} and
\texttt{FieldDecl} storage whose canonical type is a project enum or
\texttt{std::atomic<Enum>}. An assignment is a candidate only if its left-hand
side resolves through \texttt{DeclRefExpr} or \texttt{MemberExpr} to that
recorded declaration. This declaration linkage avoids choosing a write merely
because it contains a familiar spelling.

The instrumenter accepts only supported standalone writes. It captures the
right-hand expression once, performs the original assignment once, then calls
the unmodified SGFuzz callback with the captured value. Rewriting happens in a
disposable copied tree; the upstream checkout is never edited. The current
implementation supports direct and selected overloaded C++ assignments,
including the IXWebSocket \texttt{std::atomic<ReadyState>} field
\texttt{\_readyState}. It also distinguishes anonymous C enums using resolved
source coordinates.

\section{Methodology}

\subsection{Pinned inputs and evidence}

The experiments use SGFuzz commit \texttt{35f1ffc9} and its unchanged
\texttt{libsfuzzer.a} runtime. IXWebSocket is pinned at \texttt{514a0b96};
H2O is pinned at \texttt{1e7344932}. The full revision identifiers and the
recorded environment
uses Debian Clang 14.0.6, Python 3.11.2, and sanitizer-enabled C++14 builds.
The source, plans, commands, manifests, logs, and derived tables are retained
in the accompanying repository~\cite{artifact}.

\subsection{Regression controls}

Five pre-registered controls compare the regex and AST selectors using fixed
seed inputs and \texttt{-runs=0}: a direct unscoped C++ enum, a scoped C++
enum, an atomic enum, a direct C enum, and an anonymous C enum. The first and
fourth are preservation controls; the remaining cases test forms that the
regex selector misses. For each condition we record callback sites, STT nodes,
and leaves. The expected outputs were fixed before execution in
\texttt{experiments/regression-controls.json}.

\subsection{Real-project checks}

The primary real target is IXWebSocket's transport implementation. For the
deterministic feasibility pilot, control and AST builds compile the same frozen
36-unit default-library source list, fixed four-input corpus, sanitizer flags,
and SGFuzz runtime. The sole difference is that the AST condition rewrites a
copied \texttt{IXWebSocketTransport.cpp}. Both execute a deterministic
\texttt{-runs=0} replay.

A separate static audit measures parser and candidate-discovery coverage, not
fuzzing effectiveness. It parses 36 IXWebSocket C++ units with recorded
library flags and 124 H2O C units from the \texttt{h2o-fuzzer-http2} build
closure with a fresh CMake compilation database. A candidate write is an AST
assignment whose lvalue resolves to selected enum storage. It is not assumed to
be a semantic protocol state or a build-validated rewrite site.

Finally, we ran a 60-second operational preflight on the AST-instrumented
IXWebSocket transport harness using a fixed libFuzzer seed, address and
undefined-behavior sanitizers, and a 25-second per-input timeout. It records
source integrity, callback count, final STT, libFuzzer counters, corpus size,
and peak resident memory. This short campaign demonstrates that the full path
can run continuously; it is not a comparison against regex SGFuzz.

\section{Results}

\subsection{Regression controls and deterministic replay}

Table~\ref{tab:controls} reports the completed pre-registered controls. All
five passed. The direct unscoped C++ and direct C controls preserved the
baseline behavior. The AST selector recovered the scoped, atomic, and
anonymous-C cases that remained root-only under regex selection.

\begin{table}[t]
\centering
\caption{Pre-registered regression controls. Entries are callback sites / STT
nodes / leaves; node counts include the root.}
\label{tab:controls}
\footnotesize
\setlength{\tabcolsep}{2.6pt}
\begin{tabular}{@{}p{0.43\columnwidth}cc@{}}
\toprule
\textbf{Case} & \textbf{Regex} & \textbf{AST} \\
\midrule
Direct unscoped C++ control & 1 / 4 / 1 & 1 / 4 / 1 \\
Scoped C++ enum & 0 / 1 / 1 & 1 / 4 / 1 \\
Atomic C++ enum & 0 / 1 / 1 & 1 / 4 / 1 \\
Direct C enum control & 1 / 5 / 1 & 1 / 5 / 1 \\
Anonymous C enum & 0 / 1 / 1 & 2 / 3 / 1 \\
\bottomrule
\end{tabular}
\end{table}

The same pattern appears in the primary real target (Table~\ref{tab:ix}).
The regex control placed no callback in the target transport file and its
replay remained root-only. The AST condition placed five callbacks, built the
same 36-unit source scope, and produced seven nodes and two leaves under the
same four-input deterministic replay. This confirms that at least one selected
site resolves to and executes through the original SGFuzz runtime. It does not
show that all five sites execute, nor does it establish higher coverage,
throughput, or bug discovery.

\begin{table}[t]
\centering
\caption{IXWebSocket deterministic feasibility pilot. Both conditions use the
same 36-unit build, corpus, sanitizer flags, runtime, and \texttt{-runs=0}.}
\label{tab:ix}
\footnotesize
\begin{tabular}{@{}lrrrr@{}}
\toprule
\textbf{Condition} & \textbf{Callbacks} & \textbf{Nodes} & \textbf{Leaves} & \textbf{Build/replay} \\
\midrule
Control & 0 & 1 & 1 & Pass \\
AST & 5 & 7 & 2 & Pass \\
\bottomrule
\end{tabular}
\end{table}

\subsection{Static audit}

Table~\ref{tab:audit} reports the static-audit inventory. Every frozen
translation unit parsed successfully. The candidate and storage counts make
the selection scope auditable, but they are not precision labels and should
not be read as a count of true protocol states or inserted callbacks. The H2O
scan additionally shows that the analyzer can preserve per-unit CMake
arguments for a real C target rather than relying on a single guessed command.

\begin{table}[t]
\centering
\caption{Real-project AST static audit. Candidate writes are type-resolved
assignment candidates, not semantic state labels.}
\label{tab:audit}
\footnotesize
\setlength{\tabcolsep}{2.3pt}
\begin{tabular}{@{}lrrrrr@{}}
\toprule
\textbf{Scope} & \textbf{TU} & \textbf{Fail.} & \textbf{Enums} & \textbf{Storage} & \textbf{Writes} \\
\midrule
IXWebSocket C++ & 36/36 & 0 & 18 & 17 & 17 \\
H2O C & 124/124 & 0 & 98 & 73 & 171 \\
\bottomrule
\end{tabular}
\end{table}

\subsection{Short operational preflight}

The completed 60-second AST campaign returned zero, reported no AddressSanitizer
or UndefinedBehaviorSanitizer finding, and preserved the IXWebSocket source
SHA-256 before and after instrumentation. It inserted five callbacks and
finished with 2,983 STT nodes and 1,026 distinct leaves. LibFuzzer reported
\texttt{cov: 88} and \texttt{ft: 278}; these counters describe only this small
transport harness. The run executed 6,258,463 inputs in 65.007 seconds
(approximately 96k executions per second over measured elapsed time), produced
2,649 corpus files totaling 33,853 bytes, and recorded 476 MB peak RSS.

This result is operational evidence, not a randomized performance comparison.
The 23-hour single-trial plan records the same metrics at 15 minutes, one hour,
and completion; it is not treated as a result here because it had not completed
when this paper was prepared.

\subsection{Instrumentation cost and limitations}

In a separate three-repetition atomic-fixture check, median copy-and-instrument
time was 19 ms for regex and 198 ms for the AST version, a 10.4$\times$
preprocessing increase. This is not fuzzing throughput: it measures a small
one-time source preprocessing task before a campaign begins. It provides no
basis for claiming a tenfold decrease in executions per second or a target-scale
scalability result.

The AST rule is type-aware, not semantically omniscient. Some selected enum
storage can be irrelevant to useful protocol state, so over-selection may
increase callback and STT bookkeeping without improving feedback quality. The
prototype also conservatively excludes type aliases, macros, header-only
writes, atomic \texttt{.store()} calls, non-standalone assignments, and some
same-short-name enum cases. Textual-AST parsing and incomplete source locations
remain implementation boundaries. These constraints motivate manual labels,
held-out language cases, and matched multi-seed fuzzing comparisons as future
work.

\section{Conclusion}

Clang AST declaration and type linkage improves the source-selection step of
SGFuzz while preserving the original callback and STT runtime. Across five
pre-registered C and C++ controls, the prototype preserved direct-enum behavior
and recovered scoped, atomic, and anonymous-enum writes missed by regex
selection. On IXWebSocket, it inserted five target callbacks where the regex
condition inserted none, and an identical deterministic replay produced a
non-root STT. The completed short sanitizer campaign further shows that the
copied, instrumented target runs continuously without changing its upstream
source.

These findings establish improved instrumentation fidelity and end-to-end
feasibility, not superior fuzzing effectiveness. A complete comparison should
use matched corpora, seeds, CPU allocation, and trial budgets, then report STT
growth, edge coverage, executions per second, peak memory, and deduplicated
crashes with uncertainty intervals.

\section*{Individual Contributions and Attribution}

Will D'Andrade developed the project framing and AST-aware direction, organized
the report, and designed the evaluation questions. Jose Luis Gonzalez Paniagua
implemented the AST detector and copy-only instrumenter, created the regression
cases and runners, and collected the evidence. Both authors reviewed the
evidence-backed claims. Reused SGFuzz components remain attributed to the
original project~\cite{ba2022}; material AI assistance and verification are
documented in \texttt{notes/ai-use.md}.

\begin{thebibliography}{00}
\bibitem{ba2022} J. Ba, M. B\"ohme, Z. Mirzamomen, and A. Roychoudhury,
``Stateful Greybox Fuzzing,'' in \emph{Proc. 31st USENIX Security Symp.},
2022, pp. 3255--3272. [Online]. Available:
\url{https://www.usenix.org/conference/usenixsecurity22/presentation/ba}
\bibitem{clangast} Clang Project, ``Introduction to the Clang AST,''
\emph{Clang Documentation}. [Online]. Available:
\url{https://clang.llvm.org/docs/IntroductionToTheClangAST.html}
\bibitem{ixwebsocket} Machine Zone, ``IXWebSocket,'' GitHub repository.
[Online]. Available: \url{https://github.com/machinezone/IXWebSocket}
\bibitem{h2o} H2O Project, ``H2O: The Optimized HTTP Server,'' GitHub
repository. [Online]. Available: \url{https://github.com/h2o/h2o}
\bibitem{libfuzzer} LLVM Project, ``libFuzzer: A Library for Coverage-Guided
Fuzz Testing,'' \emph{LLVM Documentation}. [Online]. Available:
\url{https://llvm.org/docs/LibFuzzer.html}
\bibitem{artifact} J. L. Gonzalez Paniagua and W. D'Andrade,
``ClangAST-SGFuzz artifact,'' 2026. [Online]. Available:
\url{https://github.com/Jlg1061/clangAST-SGFuzz}
\end{thebibliography}

\end{document}
